\subsection{CSR registrai, privilegijų režimai ir išimtys}
\label{sec:csr-ir-isimtys}

Privilegijuotos architektūros realizacija sudaryta iš kelių dalių: privilegijų režimo būsenos, CSR registrų žemėlapio, pertraukimų parinkimo ir trap logikos. Palaikomi trys darbo režimai -- User, Supervisor ir Machine. Šių režimų pakanka OpenSBI, Linux branduolio ir vartotojo erdvės programų vykdymui.

CSR prieiga realizuota per bendrą registrų žemėlapį. Kiekvienam palaikomam CSR registrui priskiriamas numeris, skaitymo kelias ir rašymo kelias. Tokiu būdu \texttt{CSRRW}, \texttt{CSRRS}, \texttt{CSRRC} bei jų immediate variantai naudoja bendrą CSR skaitymo ir rašymo mechanizmą. Prieš prieigą patikrinama CSR privilegijų klasė: jeigu dabartinio režimo kodas mažesnis už CSR adrese užkoduotą reikalavimą, prieiga atmetama.

Realizacijoje palaikomi Linux paleidimui svarbiausi CSR registrai: \texttt{mstatus}, \texttt{sstatus}, \texttt{mtvec}, \texttt{stvec}, \texttt{mepc}, \texttt{sepc}, \texttt{mcause}, \texttt{scause}, \texttt{mtval}, \texttt{stval}, \texttt{medeleg}, \texttt{mideleg}, \texttt{mie}, \texttt{mip}, \texttt{sie}, \texttt{sip}, \texttt{mscratch}, \texttt{sscratch}, \texttt{satp}, \texttt{mcounteren}, \texttt{scounteren}, \texttt{time} ir eksperimentinis \texttt{qpuctx}. Kai kurie registrai, pavyzdžiui \texttt{pmpcfg0}, \texttt{pmpaddr0}, \texttt{marchid} ar \texttt{mimpid}, realizuoti kaip stub'ai, nes Linux paleidimo metu jų pakanka tik skaityti arba rašyti be realaus poveikio.

Išimčių aptarnavimas vykdomas pagal RISC-V delegavimo taisykles. Jeigu išimtis kyla ne M režime ir atitinkamas bitas pažymėtas \texttt{medeleg}, ji perduodama S režimui; kitu atveju valdymas perduodamas M režimui. Perduodant valdymą į trap handler'į, emuliatorius įrašo \texttt{mepc}/\texttt{sepc}, \texttt{mcause}/\texttt{scause}, \texttt{mtval}/\texttt{stval}, išjungia atitinkamą pertraukimų leidimą ir pakeičia dabartinį privilegijų režimą.

Pertrauktys atnaujinamos iš įrenginių būsenos. CLINT nustato mašininio laikmačio pertrauktį, kai \texttt{mtime > mtimecmp}. PLIC gali nustatyti mašininę arba supervisor išorinę pertrauktį priklausomai nuo PLIC konteksto. Pertraukčių parinkimas atliekamas pagal paprastą prioritetų sąrašą, atskirai tikrinant, ar konkreti pertrauktis turi būti aptarnauta M ar S režime.

Svarbu, kad ši realizacija yra praktinis Linux paleidimui reikalingas poaibis, o ne visos privilegijuotos specifikacijos modelis \cite{riscv_privileged_isa}. Pavyzdžiui, \texttt{WFI} vykdomas kaip \texttt{nop}, \texttt{mstatus.MPRV}, \texttt{TVM}, \texttt{TW}, \texttt{TSR}, PMP, hiperprivilegijuotas režimas ir keli naujesni būsenos laukai nėra pilnai realizuoti. Šie ribojimai priimtini šiame darbe, nes paleidžiama vieno harto Linux sistema jų nenaudoja kaip būtinos funkcijos.
